\section{Implementation Details}
\label{sec:imp}

\subsection{Architecture}

\subsubsection{Refinement Transformer}

The refinement transformer uses $I = 3$ refinement blocks, and the dimension of all embeddings is $d = 128$.
We forecast $T_f = 10$ future timestamps at stride of $s_f = 0.5\text{s}$, yielding a 5s prediction horizon.
Unless otherwise stated, we use target timestamps of $\mathcal{T}_m = \{-0.3s, -0.6s, \dots, -2.4s\}$ (i.e., $s_m = 0.3\text{s}$, $T_m = 8$) for reading from the memory bank at inference.
In the memory cross attention we use the nearest $k=4$ neighbors.
Following prior works~\cite{hednet,safdnet,centerpoint}, for any detection post-processing, we use a 0.1 confidence threshold; per-class NMS IoU thresholds of $\{0.75, 0.6, 0.55\}$ for vehicles, pedestrians, and cyclists, respectively; and a top $K=500$.
\ourmodel{} has $3.8$M parameters, while the base detectors have anywhere from $8$M (Centerpoint~\cite{centerpoint}) to $53$M (BEVMap~\cite{chang2024bevmap}) parameters.
% centerpoint: 8M
% safdnet: 10M
% hednet: 13M
% fcos3d: 53M 
% bevmap: 53M 


\subsubsection{Trajectory Forecasting Header:}

% Inspired by \cite{casas2024detra} and \cite{chitta2022transfuser},
Each refinement block predicts trajectory forecasts from the refinement features $\mathbf{Q}^{\mathrm{ref}(i)} \in \mathbb{R}^{N^\mathrm{ref} \times T_f + 1 \times d}$ using a
single-layer bi-directional GRU~\cite{cho2014learning}.
The output of the GRU is %trajectories
a sequence of 2D BEV positions for each actor of shape $\mathbb{R}^{N^\mathrm{ref} \times T_f \times 2}$,
and we use finite differences to calculate the heading.
The result is trajectories
$\mathbf{T}^{\mathrm{ref}(i)} \in \mathbb{R}^{N^\mathrm{ref} \times T_f \times 3}$,
describing each objects' BEV pose $\{(x, y, \theta)_{t + s_f}, \dots, (x, y, \theta)_{t + s_f + T_f - 1}\}$.
Note while the GRU sequence length is $T_f + 1$, we only use the last $T_f$ outputs as the future trajectory forecasts.


\subsection{Training}

\subsubsection{Scheduling}

For each base detector, we train \ourmodel{} for 60k update steps (roughly equivalent to 6 epochs on WOD and AV2), with batch size 16.
We use a cosine learning rate decay with a max learning rate of $8 \times 10^{-4}$,
and a linear warm-up for the first 1000 steps, beginning with a learning
rate of $8 \times 10^{-5}$.
During training,
we use a variable set of memory target timestamps
$T_m \sim \texttt{uniform}(\{6, 7, 8, 9, 10\})$
with a variable stride $s_m \sim \texttt{uniform}(\{0.2s, 0.3s, 0.4s\})$.

\subsubsection{Augmentations}

When training \ourmodel{} we adopt standard data augmentations used in prior works~\cite{safdnet,hednet,wang2023dsvt},
including random translations in the x, y, and z directions in the ego coordinate frame with a standard deviation of 0.5m,
random rotation around the z-axis (in the BEV plane) uniformly sampled from $[-\pi/4, \pi/4]$,
random scaling sampled uniformly from $[0.95, 1.05]$,
and flipping with $\{\text{no flip}, \text{x flip}, \text{y flip}, \text{x = y flip}\}$
occurring with equal probability.
We do not use any ground-truth sampling~\cite{yan2018second} because
it is difficult to make consistent and realistic over time and thus incompatible with memory training.

\subsection{3D Detector Implementation Details} \label{sec:off-the-shelf-imp}

We use the official implementations for all 3D detectors in our experiments:
single stage CenterPoint~\cite{centerpoint}\footnote{\href{https://github.com/tianweiy/CenterPoint}{https://github.com/tianweiy/CenterPoint}},
HEDNet~\cite{hednet} and SAFDNet~\cite{safdnet}\footnote{\href{https://github.com/zhanggang001/HEDNet}{https://github.com/zhanggang001/HEDNet}},
FCOS3D~\cite{wang2021fcos3d}\footnote{\href{https://github.com/jjw-DL/mmdetection3d\_Noted/tree/master/configs/fcos3d}{https://github.com/jjw-DL/mmdetection3d\_Noted/}},
and BEVMap~\cite{chang2024bevmap}\footnote{\href{https://github.com/mincheoree/BEVMap}{https://github.com/mincheoree/BEVMap}}.

For the LiDAR based models (CenterPoint~\cite{centerpoint}, HEDNet~\cite{hednet}, SAFDNet~\cite{safdnet}),
we train with their official configurations for the WOD.
Neither FOCS3D nor BEVMap provide official configurations for training on AV2;
we train with the ``front ring camera" on AV2 for 3 epochs with a batch size of 16, using the proposed learning rate scheduler.

To obtain $\mathbf{Q}^\mathrm{det}$ from the 3D detectors,
for each object proposal $(x, y, z, l, w, h, \theta)$ in $\mathbf{B} \in \mathbb{R}^{N \times 7}$,
we interpolate the features map directly before the detection header at the projected object centroid $(x, y, z)$.
Concretely, for detectors that produce BEV feature maps before their header
(e.g., CenterPoint~\cite{centerpoint}, HEDNet~\cite{hednet}, SAFDNet~\cite{safdnet}, BEVMap~\cite{chang2024bevmap}),
we perform bi-linear interpolation of the feature maps at the BEV object centroid $(x, y)$.
For detectors that use image-view feature maps, we perform bi-linear
interpolation at the object centroid $(x, y, z)$ projected onto the image plane.
While not experimented with in this work, we could similarly use tri-linear interpolation for detectors that use 3D feature maps.